\section*{Capítulo \ref{ch:b1:p-block-13-15} --- El bloque p I: los grupos del boro, del carbono y del nitrógeno}

\begin{solution}{exo:b1:p-block-13-15:1}
\ce{NH4NO3}: $-$III en \ce{NH4+}, $+$V en \ce{NO3-}. \ce{N2O5}: $+$V.
\ce{HNO2}: $+$III. \ce{N2H4}: $-$II. \ce{Mg3N2}: $-$III.
\end{solution}

\begin{solution}{exo:b1:p-block-13-15:2}
\ce{BF3}: boro con tres enlaces y seis electrones, trigonal plana.
\ce{NH3}: tres enlaces y un \omterm{def:b1:lewis-resonance-vsepr:lewis}{par no enlazante}. El \omterm{def:b1:lewis-resonance-vsepr:lewis}{par no enlazante} del N
forma un \omterm{def:b1:lewis-resonance-vsepr:lewis-acid}{enlace dativo} con el B: \ce{BF3} es el \omterm{def:b1:lewis-resonance-vsepr:lewis-acid}{ácido de Lewis}, y
\ce{NH3}, la \omterm{def:b1:lewis-resonance-vsepr:lewis-acid}{base de Lewis}; en el aducto, ambos átomos son tetraédricos.
\end{solution}

\begin{solution}{exo:b1:p-block-13-15:3}
\ce{CO2 + 2NaOH -> Na2CO3 + H2O}; \ce{SiO2 + 2NaOH -> Na2SiO3 + H2O};
\ce{P4O10 + 12NaOH -> 4Na3PO4 + 6H2O}; \ce{Al2O3 + 6HCl -> 2AlCl3 + 3H2O};
\ce{Al2O3 + 2NaOH + 3H2O -> 2NaAl(OH)4}.
\end{solution}

\begin{solution}{exo:b1:p-block-13-15:4}
$\Delta_r G^\circ = 2 \times (-16.45) = \qty{-32.9}{kJ/mol}$, $\log K =
32.9/5.708 = 5.76$, $K = 10^{5.76}$.
\end{solution}

\begin{solution}{exo:b1:p-block-13-15:5}
Anillo de seis unidades \ce{SiO3^2-}: \ce{Si6O18^12-}. Cadena doble, por
cada cuatro silicios: dos con dos vértices compartidos (3 O, carga $-2$
cada uno) y dos con tres ($2.5$ O, carga $-1$ cada uno): \ce{Si4O11^6-}.
\end{solution}

\begin{solution}{exo:b1:p-block-13-15:6}
\qty{550}{kg} de \ce{Al(OH)3} son $550/78.0 = \qty{7.05}{kmol}$, que dan
\qty{3.53}{kmol} de \ce{Al2O3}, \qty{360}{kg}; al \qty{92}{\percent}:
\qty{331}{kg}.
\end{solution}

\begin{solution}{exo:b1:p-block-13-15:7}
El boro tiene un orbital vacío y acepta el \omterm{def:b1:lewis-resonance-vsepr:lewis}{par no enlazante} de un ion
hidróxido tomado del agua: \ce{B(OH)3 + 2H2O <=> B(OH)4- + H3O+}. El
protón liberado procede de una molécula de agua, no del ácido bórico.
\end{solution}

\begin{solution}{exo:b1:p-block-13-15:8}
El N pasa de $-$III a $+$II, cinco electrones por cada \ce{NH3}; cada
\ce{O2} toma cuatro. $4 \times 5 = 5 \times 4$: \ce{4NH3 + 5O2 -> 4NO + 6H2O}.
\end{solution}

\begin{solution}{exo:b1:p-block-13-15:9}
\ce{NH3 + HNO3 -> NH4NO3}. Una tonelada son $10^6/80.0 = \qty{1.25e4}{mol}$;
\qty{1.25e4}{mol} de amoniaco para neutralizar y \qty{1.25e4}{mol} para
fabricar el ácido nítrico: $2.50 \times 10^4 \times 17.0 = \qty{425}{kg}$.
\end{solution}

\begin{solution}{exo:b1:p-block-13-15:10}
Los átomos del segundo \omterm{def:b1:periodicity:table}{periodo} son pequeños: sus orbitales \textit{p} se
solapan bien lateralmente y forman enlaces $\pi$ fuertes, de modo que
\ce{N#N} y \ce{O=C=O} son moléculas pequeñas y estables. Los átomos del
tercer \omterm{def:b1:periodicity:table}{periodo} son más grandes; su solapamiento $\pi$ es pobre, y
prefieren más enlaces sencillos: los tetraedros \ce{P4}, la red de
\ce{SiO4}.
\end{solution}

\begin{solution}{exo:b1:p-block-13-15:11}
$E^\circ(\ce{PbO2}/\ce{Pb^2+}) = 1.46 > E^\circ(\ce{Cl2}/\ce{Cl-}) = 1.36$~V:
\ce{PbO2 + 4H+ + 2Cl- -> Pb^2+ + Cl2 + 2H2O}. El plomo(IV) es inestable
frente al plomo(II) (par inerte); el silicio(II) no es un estado estable,
y el \ce{SiO2} no es un \omterm{def:b1:nernst:couple}{oxidante}.
\end{solution}

\begin{solution}{exo:b1:p-block-13-15:12}
$160 \times 17.0/14.0 = \qty{194}{Mt}$ de amoniaco, que contienen $194 \times
3.0/17.0 = \qty{34}{Mt}$ de hidrógeno, obtenido sobre todo del gas natural
(metano y vapor de agua).
\end{solution}

\begin{solution}{pb:b1:p-block-13-15:1}
\textbf{1.} :N\,$\equiv$\,N: --- un enlace triple, muy fuerte y sin
polaridad.
\textbf{2.} 0, $-$III, $+$II, $+$IV, $+$V, y $-$III/$+$V en \ce{NH4NO3}.
\textbf{3.} Las bacterias de las raíces de las leguminosas, los rayos y
el proceso Haber--Bosch.
\textbf{4.} Romper el enlace triple requiere un \omterm{def:b1:elementary-steps:catalyst}{catalizador} que las
plantas no tienen; las bacterias con la enzima nitrogenasa, sí.
\textbf{5.} El nitrógeno se reduce (de 0 a $-$III), y el hidrógeno se
oxida (de 0 a $+$I).
\textbf{6.} \ce{N2 + 3H2 <=> 2NH3}.
\textbf{7.} $K = 10^{2 \times 16.45/5.708} = 10^{5.76}$.
\textbf{8.} A \qty{25}{\celsius} la reacción es demasiado lenta; el
\omterm{def:b1:elementary-steps:catalyst}{catalizador} solo funciona en caliente, y las condiciones son un
compromiso que se discute en el volumen del segundo año.
\textbf{9.} \qty{58.8}{kmol} de amoniaco: \qty{29.41}{kmol} de \ce{N2},
\qty{824}{kg}; \qty{88.2}{kmol} de \ce{H2}, \qty{176}{kg}.
\textbf{10.} $V(\ce{N2}) = 29.4 \times 10^3 \times 8.314 \times 298.15/10^5 =
\qty{729}{m^3}$; aire, $729/0.78 = \qty{935}{m^3}$.
\textbf{11.} Del gas natural (reformado del metano con vapor), con
\ce{CO2} como subproducto.
\textbf{12.} Inyectado en el suelo es un fertilizante barato, pero es un
gas tóxico; los sólidos (urea, nitrato de amonio) son más fáciles de
almacenar y de esparcir.
\textbf{13.} N: $-$III $\to$ $+$II (5 electrones); O: 0 $\to$ $-$II (4 por
\ce{O2}): \ce{4NH3 + 5O2 -> 4NO + 6H2O}.
\textbf{14.} \ce{2NO + O2 -> 2NO2}.
\textbf{15.} \ce{3NO2 + H2O -> 2HNO3 + NO}: una \omterm{def:b1:e-ph-diagrams:disproportionation}{dismutación} (N de $+$IV a
$+$V y a $+$II).
\textbf{16.} \ce{NH3 + 2O2 -> HNO3 + H2O}.
\textbf{17.} $2 \times 58.8 \times 32.0 = \qty{3.76}{t}$.
\textbf{18.} Hace la oxidación rápida y selectiva hacia el NO, en lugar
del \ce{N2}, más estable.
\textbf{19.} \ce{NH3 + HNO3 -> NH4NO3}.
\textbf{20.} La mitad.
\textbf{21.} \qty{29.4}{kmol} de \ce{NH4NO3}, \qty{2.35}{t}.
\textbf{22.} $28.0/80.0 = \qty{35}{\percent}$.
\textbf{23.} Contiene un \omterm{def:b1:nernst:couple}{oxidante} (el nitrato) y un \omterm{def:b1:nernst:couple}{reductor} (el amonio)
en el mismo \omterm{def:b1:crystals-metals:crystal}{cristal}: calentado o mezclado con combustibles, puede
descomponerse de forma explosiva.
\textbf{24.} $58.8 \times 63.0 =$ \textbf{\qty{3.7}{t} de ácido nítrico por
tonelada de amoniaco}.
\end{solution}
